% Generated draft: requires mathematical and source review.
\documentclass[11pt]{article}
\usepackage[margin=25mm]{geometry}
\usepackage{fontspec}
\setmainfont{Latin Modern Roman}
\usepackage{amsmath,amssymb,mathtools,unicode-math}
\setmathfont{Latin Modern Math}
\usepackage{enumitem,needspace,xurl,hyperref}
\setlength{\parindent}{0pt}
\setlength{\parskip}{5pt}
\setlength{\emergencystretch}{3em}
\newcommand{\N}{\mathbb{N}}
\newcommand{\Z}{\mathbb{Z}}
\newcommand{\Q}{\mathbb{Q}}
\newcommand{\R}{\mathbb{R}}
\newcommand{\C}{\mathbb{C}}
\newcommand{\F}{\mathbb{F}}
\newcommand{\A}{\mathbb{A}}
\newcommand{\PP}{\mathbb{P}}
\newcommand{\E}{\mathbb{E}}
\newcommand{\eps}{\varepsilon}
\newcommand{\dd}{\,\mathrm{d}}
\begin{document}
\section*{UnsolvedMath ID: 2200020}
\textbf{Source problem code:} catalogue entry
\subsection{Definitions and mathematical statement}
Let $p\in\R[x]$ be a nonzero polynomial of degree $k$. The source associates polynomials $P_i(x)$ to $p$ and conjectures
\[
\#_r P_i\ \le\ \min\{\deg P_i,k\}
\]
for every index $i$ in the intended family. Its supplied summary calls these Jensen polynomials, but the record contains no formula defining $P_i$, no index range, and no normalization. The symbol cannot be identified with some other standard construction bearing the name Jensen polynomial without further source information.

The real-zero counting convention is also absent. For a nonzero polynomial $Q$, two possible meanings are
\[
\#_r^{\mathrm{dist}}Q=\#\{x\in\R:Q(x)=0\},\qquad
\#_r^{\mathrm{mult}}Q=\sum_{\substack{x\in\R\\Q(x)=0}}\operatorname{ord}_x(Q).
\]
The displayed conjecture retains the source symbol $\#_r$ pending a choice; it does not silently replace a multiplicity count by a distinct-root count. Both the degree and ordinary root counts require separate conventions for an identically zero $P_i$, and the source does not say which indices avoid this case. The admissible degrees $k$, including constant $p$, also need confirmation.

The supplied research summary labels the full $P_i$ assertion open. It reports that the $i=1$ inequality $\#_rP_1\le\#_{nr}p$, together with the stated bounds $\deg P_1\le2k-2$ and $\#_{nr}p\le k$, is a proved Hawaii-conjecture base case attributed to Tyaglov (2011). Here $\#_{nr}p$ denotes the intended number of nonreal zeros, again with an unspecified multiplicity convention. The exact definitions and hypotheses must be checked before this can be used as a theorem about the missing $P_1$.

The summary separately mentions a conjectural adjoining family
\[
\#_rG_i\le\min\{\deg G_i,\#_{nr}p\},
\]
called Corollary 1 in the source, and reports a 2024 counterexample to its $i=1$ member by Katkova--Tyaglov--Vishnyakova. The $G_i$ are also undefined here. That reported counterexample is not asserted to refute the distinct $P_i$ conjecture. All literature-status statements are supplied assessments, not independent verification.
\subsection{Short English statement}
For the source's family $P_i$ associated to a degree-$k$ real polynomial, bound the number of real zeros by both its degree and $k$. The definitions of that family, its indices, and root multiplicities are missing, so the assertion remains conditional. A reported counterexample for a neighboring $G_i$ family must be kept separate from this conjecture.
\subsection{Sources}
\begin{itemize}
\item Local source registry: \url{data/unsolvedmath/problems.json}, permanent UnsolvedMath ID $2200020$, source code catalogue entry, including the supplied OPEN-TRIAGE summary. Dataset attribution: Ulam AI and UnsolvedMath Contributors; original-source attribution and rights are retained.
\item Shapiro, \emph{Problems Around Polynomials: The Good, The Bad and The Ugly} (2015), source-order item 20, Conjecture 13 as printed, attributed in the extract to B. Sh: \url{https://arxiv.org/abs/1503.05295}. The metadata reports source-TeX extraction and July 29, 2026 access, with status and rights review flags.
\item Tyaglov (2011), cited by the supplied summary for the Hawaii conjecture and the stated $i=1$ base case. No full title, theorem statement, or separate link is supplied.
\item Katkova--Tyaglov--Vishnyakova, arXiv:2406.00686 (2024), cited by the summary for a counterexample to the adjoining $G_1$ claim. The record gives no full title or definitions of that family. No linked source was inspected for this draft.
\end{itemize}
\end{document}
